\documentclass[11pt,a4paper]{article}
\usepackage[margin=2.2cm]{geometry}
\usepackage[T1]{fontenc}
\usepackage{amsmath,amssymb,xcolor,enumitem,booktabs,graphicx,tabularx,array}
\usepackage[normalem]{ulem}
\graphicspath{{../../../../thesis/Images/}}
\setlength{\parskip}{0.6em}\setlength{\parindent}{0pt}
\definecolor{newc}{HTML}{1f6f43}
\newcommand{\new}[1]{\textcolor{newc}{#1}}
\newcommand{\old}[1]{\textcolor{gray}{\sout{#1}}}
\newcommand{\note}[1]{{\small\itshape [#1]}}
\newcommand{\lab}[1]{\textsf{\footnotesize\detokenize{#1}}}
\renewcommand{\ref}[1]{\lab{#1}}
\renewcommand{\eqref}[1]{(\lab{#1})}
\newcommand{\nref}[1]{\lab{#1}}
\newcommand{\citep}[1]{[\lab{#1}]}
\newcommand{\citet}[1]{\lab{#1}}
\newcommand{\where}[1]{\par\textbf{#1}\par}
\title{Proposed change, Problem Formulation, round 3\\ Seams after the cuts, and figure sizes}
\date{25 November 2026, version 1}
\begin{document}\sloppy
\maketitle

\section*{How to read this}
I read the compiled chapter after [A133] and [A134], and the Method
paragraphs that received text. Part~1 fixes six seams the cuts left. Part~2
proposes smaller figures. Black text stays, \new{green} is new and
\old{struck} text goes. \note{Notes in brackets} say why. Answer by number.
Nothing below is in the thesis yet.

\textbf{What it saves, honestly.} Part~1 saves a few lines and is about
reading, not length. Part~2 saves about 9\,cm of figure height, a little
under half a page, because the chapter is mostly text. Its main gain is that
the pages look lighter. To save a full page more, a figure would have to go
(G10).

% ================================================================
\section*{Part 1. Seams}

\subsection*{S1. $\mathcal{F}$ defined once (Fix)}
\where{Where the items are, after the three conditions}
\dots\ so a line that waited for them waits for a later one. \old{A storage
rule is any $F$ that meets these three conditions, and $\mathcal{F}$ is the
set of them.} \new{A placement that meets these three conditions is a valid
one.}

\where{Orders and tasks, unchanged}
A storage rule therefore does two things. It places arriving stock at a
storage update, and it ranks the shelves a line takes from. From here on $F$
means both parts, and $\mathcal{F}$ holds every rule that places stock under
the three conditions of a storage update and ranks the shelves for every
line.

\note{The reader met two definitions of $\mathcal{F}$, four pages apart.
Now the storage update says what a valid placement is, and $\mathcal{F}$ is
defined once, where both parts are known.}

\subsection*{S2. The order estimate next to its words (Fix)}
\where{What the rule reads}
The first is how often each SKU was ordered. Write $\hat\rho_t(k)$ for the
order lines of SKU $k$ per step in the last storage period. If tea was
ordered 30 times in a period of 100 steps, $\hat\rho_t(\mathrm{tea})=0.3$.
\new{Formally, the order estimate is}
\[ \new{\hat\rho_t(k) = \frac{\#\{\text{order lines for } k \text{ in the last period}\}}{\Delta} .} \]

The second is how long agents stood still at each shelf. \dots\ so the time
its detour takes is not counted.

Demand storage reads only the first number. Congestion-aware storage reads
both. \old{Formally, the order estimate is \dots\ \eqref{eq:estimates}}

\note{After R7 the formula hung after a two-sentence paragraph about both
numbers, and it only defines the first. It moves to where $\hat\rho$ is
explained. The label stays.}

\subsection*{S3. The field of view's limit in one place (Fix)}
\where{The field of view}
The agent sees every vertex it can reach within a travel cost of
$d_{\mathrm{obs}}$, and the edges between them. So it knows which moves are
open nearby, not only which vertices exist. It sees nothing further away\new{,
so it cannot see a conflict coming from further than $d_{\mathrm{obs}}$}.

\where{After \lab{fig:observation}}
\old{The agent cannot see a conflict coming from further away than
$d_{\mathrm{obs}}$.}

\note{After R3 this was a one-sentence paragraph under the figure.}

\subsection*{S4. ``So'' that does not follow (Fix)}
\where{A task's life}
Only a free agent, one with no active task, can get a new task\old{. So
agent $a_2$ is free again once $\tau_1$ is delivered.}\new{, so $a_2$ can take
its next task once $\tau_1$ is delivered.}

\note{Being free does not follow from the first sentence. Taking the next
task does.}

\subsection*{S5. Method, crowding in its own paragraph (Fix)}
\where{How an agent chooses its move, ``What the network reads''}
After R6 one paragraph runs vertex features, then crowding, then the
earlier agents' vertices, which are vertex features again. The crowding
sentences move to a paragraph of their own after it:

\dots\ Together they are the \emph{vertex features} $f(v,t)$. \old{The
crowding feature \dots\ zone edge.} The vertices chosen by earlier agents
\dots\ so every controller reads the same list of features.

\new{\textbf{Crowding.} The crowding feature $\delta_i(t)$ is added once for
the agent, not for each vertex. With it a controller can route around a
crowd before an agent gets stuck in it, the argument
\citet{pham2024optimizingcrowdawaremultiagentpath} make for their
crowd-density feature. Every controller gets it, so it is measured the same
way for every agent. A zone or the centralised controller sees more agents,
so their values could be worked out, but no controller reads them.
Appendix~\ref{app:controllers} explains why a section-based agent's field of
view is not cut at the zone edge.}

\subsection*{S6. Method, the smoothing in order (Fix)}
\where{Where arriving stock goes, before \lab{eq:standstill}}
\new{A shelf sees only a few picks in one period, so one reading jumps up and
down by chance.} $\hat\omega_t(v)$ is the plain average of the shelf's last
five readings. Five is a practical smoothing choice. \old{A shelf sees only a
few picks in one period, so one reading jumps up and down by chance.} It
evens out one noisy period and still follows recent congestion, and the
number was not tuned. \dots

\note{The reason now comes before the choice, as your style asks.}

% ================================================================
\section*{Part 2. Figure sizes}
The thesis text is 15\,cm wide. Each figure below is printed at the size it
would have in the thesis, so you can judge whether it still reads. Heights
are of the image alone. Nothing is redrawn, only the width changes.

\begin{tabular}{@{}lllrr@{}}
\toprule
Item & Figure & Width now $\to$ proposed & Height now & Saved \\
\midrule
G1 & \lab{fig:instance} & 0.8 $\to$ 0.6 & 8.9\,cm & 2.2\,cm \\
G2 & \lab{fig:observation} & 0.9 $\to$ 0.7 & 8.8\,cm & 2.0\,cm \\
G3 & \lab{fig:storageupdate} & 0.8 $\to$ 0.6 & 6.1\,cm & 1.5\,cm \\
G4 & \lab{fig:pickuprule} & 0.75 $\to$ 0.6 & 5.6\,cm & 1.1\,cm \\
G5 & \lab{fig:crowding} & 0.55 $\to$ 0.45 & 5.1\,cm & 0.9\,cm \\
G6 & \lab{fig:loop} & 0.9 $\to$ 0.8 & 6.8\,cm & 0.8\,cm \\
G7 & \lab{fig:storageexample} & 0.8 $\to$ 0.7 & 6.0\,cm & 0.7\,cm \\
G8 & \lab{fig:lifecycle} & 1.0 $\to$ 0.85 & 4.7\,cm & 0.7\,cm \\
G9 & \lab{fig:conflicts} & 1.0 $\to$ 0.85 & 3.9\,cm & 0.6\,cm \\
\midrule
 & & & total & 10.5\,cm \\
\bottomrule
\end{tabular}

Kept as they are: \lab{fig:controllers-example} (three panels, already
flat), \lab{fig:taskroute} (already wrapped beside the text).

\newcommand{\gfig}[3]{\par\textbf{#1}\par\centering\includegraphics[width=#3]{#2}\par\raggedright\vspace{0.5em}}

\gfig{G1, the warehouse instance, at 9\,cm}{instance_simple.png}{9cm}
\note{The first picture of the chapter, and it has a legend. At 9\,cm the
legend is still readable. 0.65 if you want it a little larger.}

\gfig{G2, the observation, at 10.5\,cm}{observation.png}{10.5cm}
\note{The distance labels are small. Check that the blue numbers still read.}

\gfig{G3, one storage update, at 9\,cm}{storage_update.png}{9cm}

\gfig{G4, one line split into two picks, at 9\,cm}{pickup_rule.png}{9cm}

\gfig{G5, the crowding feature, at 6.75\,cm}{crowding.png}{6.75cm}

\gfig{G6, the loop a period later, at 12\,cm}{shelf_loop.png}{12cm}

\gfig{G7, arriving tea under fixed and demand storage, at 10.5\,cm}{arrival_storage.png}{10.5cm}

\gfig{G8, the life of the tea task, at 12.75\,cm}{lifecycle.png}{12.75cm}

\gfig{G9, the two conflicts, at 12.75\,cm}{conflicts.png}{12.75cm}

\subsection*{G10. If you want a page more (Choice, not recommended)}
Only removing a figure saves much more. The two candidates are the ones
whose text now says what they show.
\begin{itemize}[nosep]
  \item \lab{fig:storageupdate}. Its sentence (N4) says the whole update. It
    is the only picture of a lost pallet.
  \item \lab{fig:crowding}. \lab{fig:observation} already shows a field of
    view with agents in it. The loop section and the mean crowding refer to
    \lab{fig:crowding}, so both would need a new example.
\end{itemize}
I would keep both, since each shows something the reader otherwise only
reads about.

% ================================================================
\section*{When approved}
Applied under [A135], with a GAPS.md entry. Only widths change in the
\texttt{\textbackslash myfig} calls, so the figure scripts are not rerun. The
chapter is rebuilt, each changed page is looked at for a figure that drifted
from its text, and the tags go at the end of a line.

\end{document}
